\def\ChapterCode{MP2}
\chapter
[\CO{[MP2]} Spezielle Medizinprodukte]
{Spezielle Medizinprodukte}
\label{chp:MP2}

\ChapterIntro
{%
Das folgende Kapitel soll einen Überblick über die wichtigsten Funktionen der
entsprechenden Geräte bieten. Es ersetzt keine vom Hersteller herausgegebene
Gebrauchsanweisung. Diese ist vom Betreiber dem Personal in geeigneter Form
zugänglich zu machen und muss von den Anwendern konsultiert werden.
}
{%
\begin{description}
  \item[Maintainer:] Sebastian Gabriel
  \item[Autoren:] Michael Neuhold
  \item[Reviewer:] Standard-Reviewprozess
  \item[Version:] {\DocVersion} ({\DocVersionInfo})
  \item[SHA1:] {\DocGitCommit}
\end{description}

{\TxTeilkapitelAASS}
}


\TextSlide{\TxQuerverweise}
{~

\noindent\includegraphics[angle=270,width=\linewidth]{./images/leitsystem/Leitschema-PAM_MP_.pdf}
}{\begin{description}
  \item[Beschreibung der Medizinprodukte:] Kapitel \CO{[MP1,
  MP2]} (\ref{chp:MP1}, \ref{chp:MP2}).
  \item[Untersuchungen:] Kapitel \CO{[BAU]} (\ref{chp:BAU}).
  \item[Sanitätshilfemaßnahmen:]  Kapitel \CO{[SAN]}
  (\ref{chp:SAN}).
%   \item[Anamnese und Patientengespräch:]  Kapitel [AUP] (\ref{chp:AUP}).
\end{description}}{}{}{}{}

\SectionUnderConstruction

\section{Multifunktionsmonitore und Defibrillatoren}

% MP-MEDTRONIC-LP, % 
% MP-MEDTRONIC-LP-12, % 
% MP-MEDTRONIC-LP-15, % 
% MP-MEDTRONIC-LP-500, %
% MP-MEDTRONIC-LP-1000, %  
% , % 
% , % 
% , %  
% MP-SCHILLER-FRED, %  
% %% Beatmungsgeräte
% MP-DRAEGER-OXYLOG-1000, % 
% MP-DRAEGER-OXYLOG-2000, % 
% MP-DRAEGER-OXYLOG-3000, % 
% MP-WEINMANN-MEDUMAT-COMPACT, %
% MP-WEINMANN-MEDUMAT-STANDARD, %
% MP-WEINMANN-MEDUMAT-STANDARDA, %
% MP-WEINMANN-MEDUMAT-TRANSPORT % 

\subsection{Fred{\protect\texttrademark}}

\TakeNotesLarge{}

\opt{MP-SCHILLER-DEFIGARD}{
\subsection{Defigard{\protect\texttrademark}}

\opt{MP-SCHILLER-DEFIGARD-1002}{
\subsubsection{Defigard{\protect\texttrademark} 1002}

\TakeNotesLarge{}
}

\opt{MP-SCHILLER-DEFIGARD-2002}{
\subsubsection{Defigard{\protect\texttrademark} 2002}

\TakeNotesLarge{}
}

\opt{MP-SCHILLER-DEFIGARD-6002}{
\subsubsection{Defigard 6002}

\TakeNotesLarge{}
}

}

\subsection{Corpuls\texttrademark}



\begin{table}[H]
% \Captionof{table}{Corpuls 3: Ausstattungsoptionen.}
\begin{tabulary}{\linewidth}{LL}
\TabHeading{Merkmal}
&
\TabHeading{Beschreibung}
\\\addlinespace\midrule\addlinespace
\TabSub{Hersteller}
&
GS Elektromedizinische Geräte G. Stemple GmbH
\\\addlinespace
\TabSub{Homepage}
&
\url{http://www.corpuls.com}
\\\addlinespace\bottomrule
\end{tabulary}
\end{table}

\subsubsection{Corpuls{\protect\texttrademark} 3}

\TakeNotesLarge

Der Corpuls 3 ist ein Defibrillator/Monitoring-System. Er verlässt den
traditionellen Weg eines klassischen Kompaktgerätes und ist stattdessen modular aufgebaut.

Das System ist teilbar in:

\begin{enumerate}
    \item Monitoreinheit
    \item Patientenbox mit Messmodulen
    \item Defibrillator/Schrittmacher
\end{enumerate}


Die Patientenbox ist mit vielen Messoptionen erhältlich. Die Anschlüsse
befinden sich auf der linken oder rechten Seite der Patientenbox.

\Layout{47}{}
{}
{}
{%
\begin{tabulary}{\linewidth}{LL}
\TabHeading{Ausstattung}
&
\TabHeading{Beschreibung}
\\\addlinespace\midrule\addlinespace
\TabSub{CompactFlash} (Basisausstattung) 
&
Datenspeicherung auf
CompactFlash-Karte, Übergabe der Einsatzdaten u. a. über CompactFlash-Karte möglich
\\\addlinespace
\TabSub{12-Kanal-EKG} (Basisausstattung)
&
Simultanes 12-Kanal EKG, optional
auch mit Vermessung und Interpretationssoftware
\\\addlinespace
\TabSub{SpO2}
&
Messung von Sauerstoffsättigung durch Masimo SET-Technologie
\\\addlinespace
\TabSub{CO2}
&
Kapnometrie, Verfahren auch bei nicht-intubierten Patienten
möglich
\\\addlinespace
\TabSub{Nicht-invasive Blutdruckmessung} (NIBD)
&
Automatisierte
Messung des Blutdrucks bei Erwachsenen und Neonaten
\\\addlinespace
\TabSub{2-Kanal Temperatur}  
&
Kerntemperaturmessung invasiv und Hautoberfläche
\\\addlinespace
\TabSub{Invasive Druckmessung}
&
Bis zu 4 Kanäle IBP, gleichzeitige Messung
des arteriellen, venösen und Gehirndrucks möglich
\\\addlinespace
\TabSub{Weitere Optionen}
&
Option EKG-Vermessung und EKG-Interpretation, Option Datenübertragung mit
GSM/GPRS \\\addlinespace\bottomrule
\end{tabulary}
}
{Corpuls{\protect\texttrademark} 3: Ausstattungsoptionen}
{}







% \opt{ORG-ASBFLO}{\noindent\ASBFLO{} Es werden folgende Ausstattungsmerkmale eingesetzt:
% 
% \begin{compactenum}
%     \item CompactFlash
%     \item 12-Kanal-EKG
%     \item Sp\ce{O2}
%     \item \ce{CO2}
%     \item NIBD
% \end{compactenum}
% }


\TakeNotesLarge



\subsection{LifePak{\protect\texttrademark} AED}

\TakeNotesLarge

\subsection{LifePak{\protect\texttrademark} Multifunktionsgeräte}

\subsubsection{LifePak{\protect\texttrademark} 12}

\TakeNotesLarge{}

\opt{STATUS-EXPERIMENTAL}{
\subsubsection{LifePak{\protect\texttrademark} 15}

\TakeNotesLarge{}
}

\section{Beatmunggeräte}

\subsection{Weinmann Medumat{\protect\texttrademark}}

\begin{table}[H]
% \Captionof{table}{Corpuls 3: Ausstattungsoptionen.}
\begin{tabulary}{\linewidth}{LL}
\TabHeading{Merkmal}
&
\TabHeading{Beschreibung}
\\\addlinespace\midrule\addlinespace
\TabSub{Hersteller}
&
Weinmann Geräte für Medizin GmbH+Co.KG
\\\addlinespace
\TabSub{Homepage}
&
\url{http://www.weinmann.de}
\\\addlinespace\bottomrule
\end{tabulary}
\end{table}

\subsubsection{Medumat
Compact{\protect\texttrademark}}

\subsubsection{Medumat Standard{\protect\texttrademark} und
Medumat Standard~a{\protect\texttrademark}}

\begin{figure}[H]
\includegraphics[width=\linewidth]{./images/geraete/ger-medumat-standard-bedienfeld.jpg}
\CaptionofDiff{figure}{Bedienfelder Berieselungseinheit \E{Modul Oxygen} (li.)
und Beatmungsgerät \E{Medumat{\texttrademark} Standard} (re.).

\E{Modul Oxygen (li.):} Flowmeter, Ein-Ausschalter und Regelventil. Der
Anschluß links dient der Verbindung der Einheit mit dem Sauerstoffnetz des Fahzeuges.

\E{Medumat{\texttrademark} Standard (re.):} \Z{Air-Mix-Schalter,
Ein-Ausschalter, Stellräder für Atemfrequenz, Atemminutenvolumen (!) und Druckbegrenzung,
Beatmungsdruckanzeige, Warnleuchten für Verlegung (\enquote{Stenosis}), Lösung des
Beatmungsschlauches (\enquote{Disconnection}), niederen Flaschendruck (\enquote{<\,2.7\Bar*})
und schwache Batterie.} }
{Bedienfelder Berieselungseinheit \E{Modul Oxygen}
(li.) und Notfallbeatmungsgerät \E{Medumat{\texttrademark} Standard} (re.).}
\end{figure}

\TakeNotesLarge{}

\subsubsection{Medumat Transport{\protect\texttrademark}}

\TakeNotesLarge{}





\opt{STATUS-EXPERIMENTAL}{
\subsection{Draeger{\protect\texttrademark}
Oxylog{\protect\texttrademark}}

\subsubsection{Oxylog{\protect\texttrademark} 1000}

\TakeNotesLarge{}

\subsubsection{Oxylog{\protect\texttrademark} 2000}

\TakeNotesLarge{}

\subsubsection{Oxylog{\protect\texttrademark} 3000}

\TakeNotesLarge{}
}
\ChapterEnd